\setcounter{section}{4}
\section{Conduct a human review}
\label{sec:human-review}

\noindent\textbf{Review record:} Team review (Fredrik, Johan, Elling and Mats), 26 November 2026,
examining unified Part~2 A--D runs (\texttt{affae3e137cbe226}, \texttt{e8c0c6a8f4071325},
\texttt{9e25203dcd939450}, \texttt{unified\_part2/setting\_D}), Part~3 sweeps and diagnostics
(\texttt{054b0e6b63381969}, \texttt{3142fac3111db314}, \texttt{exec\_diagnostic}) and Part~4 replays.

\subsection*{Strongest supporting evidence}
The strongest evidence for preferring PPO in uncertain simulation is the matched Setting~B
comparison executed in one consistent computing environment: baseline PPO achieved 70.2\%
mean success against Q-learning's 31.3\% across five training seeds and 200 frozen episodes
per model, outperforming Q on all five paired seeds. In Part~3's locked confirmation on the
same 200 test seeds, PPO12 ($2\times$ learning rate, 12 epochs) raised mean success further
to 90.2\% while halving between-seed SD from 16.1 to 8.0 percentage points. Because Q and PPO
use different state encodings from the shared observation vector, this supports the tested
workflows rather than PPO in isolation.

The paired Setting~D study provides the cleanest control of observation sensitivity: the exact
same frozen Setting~B checkpoints ran on 200 paired seeds under clean and 32$\times$32 obstructed
observations with native dynamics untouched. Success fell from 70.2\% to 47.5\% for baseline PPO,
from 91.1\% to 70.1\% for PPO12, and from 27.1\% to 22.2\% for Q. Higher retained success and
smaller absolute degradation are distinct criteria, and opposite individual-seed shifts (e.g.,
PPO seed~2718 gaining $+8.0$ pp while seed~314 lost $-37.0$ pp) prevent universal claims.

\subsection*{A claim that goes beyond the evidence}
A tempting overstatement from the 90.2\% PPO12 score is ``PPO12 generalises reliably to
unfamiliar operating conditions.'' That claim exceeds our evidence: even under a single
fixed 32$\times$32 object-feature obstruction with intact maze geometry, PPO12 dropped 21.0
percentage points and failed 29.9\% of episodes (and baseline PPO failed 52.5\%). A defensible
statement is: the tested PPO workflows retain higher mean success than Q under this fixed local
obstruction, without establishing pixel-perception robustness, unseen-map transfer or deployment safety.

\subsection*{A competing explanation}
Two mechanisms required distinguishing competing explanations. First, PPO12's early lead at
65,536 interactions could have reflected noise from 10 development episodes; our 50-episode
fresh diagnostic (\texttt{3142fac3111db314}) ruled out pure episode noise (29.6\% for 12-epoch
PPO vs.\ 5.6\% for Q and 0\% for 6-epoch PPO), showing rollout reuse drives the early reversal.
Second, zero-entropy PPO's cap-level gain (64\% vs.\ 36\% on development) could reflect argmax
concentration rather than a uniformly superior policy distribution. On 50 fresh paired episodes,
the zero-entropy advantage narrowed from $+20.0$ pp under argmax (66.4\% vs.\ 46.4\%) to
$+14.0$ pp under sampling (57.2\% vs.\ 43.2\%), with lower action entropy (0.426 vs.\ 0.624 nats).
Thus execution mode accounts for part—but not all—of the cap-level gain, while seed~2718 still
reversed under argmax.

\subsection*{Comparison with submitted predictions}
Original Part~1 predictions remain unchanged with dated annotations attached. The predicted early
Q advantage held against 6-epoch PPO in B but was reversed by 12-epoch PPO and in A. Larger Q
seed variation held in B but not in A or C. Stochasticity lowered mean PPO success ($-9.8$ pp)
while leaving Q's mean unchanged (masking paired swings from $-73$ to $+44$ pp), and the common-pool
ghost probe showed only a 1.9 pp A-to-B PPO shift, leaving safer routes unestablished. Setting~C
shortened wins but reduced mean success and widened PPO SD. Same-state probes failed to support
State~1, could not verify State~2's precise-timer premise, and directly contradicted State~3
(PPO chose RIGHT 98.7\%, not LEFT). Accordingly, our initial preference for PPO becomes strictly
conditional on simulation use.

\subsection*{Initial priority rankings and group negotiation}
Table~\ref{tab:priorities} records each member's initial ranking of failure avoidance (F),
robustness to change (R), explainability (E) and time efficiency (T).

\begin{table}[H]\centering\small
\begin{tabular}{ll}\toprule
Member & Initial priority order\\\midrule
Fredrik & $F>R>E>T$\\
Johan   & $R>F>T>E$\\
Elling  & $F>E>R>T$\\
Mats    & $F>R>T>E$\\
\bottomrule\end{tabular}
\caption{Initial priority rankings reported by each team member before negotiation.}
\label{tab:priorities}
\end{table}

The negotiated final order was $F>R>E>T$. Three members placed failure avoidance first; Johan
initially ranked robustness first because of incomplete observations, weighing the paired D
degradation heavily ($-22.7$ pp for baseline PPO, $-21.0$ pp for PPO12). Keeping $F$ and $R$
above $E$ and $T$ reflected the group's consensus on avoiding failures under change. Fredrik and
Elling ranked explainability above time efficiency, whereas Johan and Mats ranked time efficiency
higher. We broke this tie using failure avoidance: in Part~4, 11 of 12 collision alternatives
still lost a life within 64 decisions (e.g., SF2/491 and TF2/756), and only TF2/941 avoided loss.
Being able to inspect whether an intervention actually prevents failure mattered more than minor
runtime savings, placing $E$ third and $T$ fourth. This ordering anchors Part~6: PPO12 is best for
further B simulation, but neither workflow is adequately supported for consequential autonomous
operation or unfamiliar environments.
